% ==============================================================
% Section 6: The KL Divergence Connection
% ==============================================================

\section{The KL Divergence Connection}
\label{sec:kl-connection}

The budget shares $w_i^*$ that emerge from optimization are not arbitrary---they are the \emph{induced weights} of the price index $\Pstar$. This observation connects our results to information theory. The Kullback-Leibler divergence between budget shares and primitive weights measures price distortion and governs the welfare value of substitutability.

% --------------------------------------------------------------
\subsection{Budget Shares as Induced Weights}
\label{subsec:induced-weights}
% --------------------------------------------------------------

Recall from the chapter on generalized means that every classical mean induces a set of weights. For a classical mean $\Mbar_{\tilde{\rho}}(\mathbf{p}; \bom)$ with power $\tilde{\rho}$, the induced weights are:
\begin{equation}
\tilde{w}_i = \frac{\omega_i \, p_i^{\tilde{\rho}}}{\sum_{j=1}^{n} \omega_j \, p_j^{\tilde{\rho}}}.
\end{equation}

The price index $\Pstar$ is a classical mean with power $\tilde{\rho} = 1 - \varepsilon = \rho/(\rho-1)$. Its induced weights are:
\begin{equation}
\tilde{w}_i = \frac{\omega_i \, p_i^{1-\varepsilon}}{\sum_{j=1}^{n} \omega_j \, p_j^{1-\varepsilon}}.
\end{equation}

Comparing with the budget share formula from Section~\ref{sec:solution}:

\begin{proposition}[Budget Shares Are Induced Weights]
\label{prop:shares-induced}
The optimal budget shares equal the induced weights in the price index:
\begin{equation}
w_i^* = \tilde{w}_i = \frac{\omega_i \, p_i^{1-\varepsilon}}{\sum_{j=1}^{n} \omega_j \, p_j^{1-\varepsilon}}.
\end{equation}
\end{proposition}

This is not a coincidence. The mathematics of generalized means and the economics of constrained optimization are two perspectives on the same structure.

% --------------------------------------------------------------
\subsection{The Zero-Distortion Benchmark}
\label{subsec:zero-distortion}
% --------------------------------------------------------------

When do budget shares equal the primitive weights? The answer characterizes a natural benchmark.

\begin{proposition}[Zero-Distortion Characterization]
\label{prop:zero-distortion}
The following conditions are equivalent:
\begin{enumerate}[label=(\roman*), leftmargin=2em]
    \item Budget shares equal weights: $w_i^* = \omega_i$ for all $i$
    \item Prices are equal: $p_1 = p_2 = \cdots = p_n$
    \item Optimal quantities equal weights: $x_i^* = \omega_i \cdot (W/p)$ where $p$ is the common price
    \item The price index equals the common price: $\Pstar = p$
\end{enumerate}
\end{proposition}

\begin{proof}
From the budget share formula, $w_i^* = \omega_i$ requires:
\[
\frac{\omega_i \, p_i^{1-\varepsilon}}{\sum_j \omega_j \, p_j^{1-\varepsilon}} = \omega_i \quad \Longleftrightarrow \quad p_i^{1-\varepsilon} = \sum_j \omega_j \, p_j^{1-\varepsilon} \quad \text{for all } i.
\]
This holds for all $i$ only if $p_i^{1-\varepsilon}$ is constant across $i$, i.e., $p_i = p$ for all $i$. The remaining equivalences follow immediately.
\end{proof}

\paragraph{Interpretation.}
Equal prices represent the ``undistorted'' benchmark. The consumer allocates according to primitive preferences $\omega_i$, with no substitution induced by price heterogeneity. Any deviation of $w_i^*$ from $\omega_i$ reflects the distorting effect of unequal prices.

% --------------------------------------------------------------
\subsection{Measuring Distortion: The KL Divergence}
\label{subsec:kl-divergence}
% --------------------------------------------------------------

The natural measure of how far budget shares deviate from primitive weights is the Kullback-Leibler divergence.

\begin{definition}[KL Divergence of Budget Shares]
\label{def:kl-shares}
The \emph{KL divergence} from primitive weights to budget shares is:
\begin{equation}
\KL{w^*}{\omega} = \sum_{i=1}^{n} w_i^* \ln \frac{w_i^*}{\omega_i}.
\end{equation}
\end{definition}

From the properties established in the generalized means chapter:

\begin{proposition}[Properties of Budget Share Divergence]
\label{prop:kl-properties}
The KL divergence $\KL{w^*}{\omega}$ satisfies:
\begin{enumerate}[label=(\roman*), leftmargin=2em]
    \item \textbf{Non-negativity}: $\KL{w^*}{\omega} \geq 0$
    \item \textbf{Zero at benchmark}: $\KL{w^*}{\omega} = 0$ if and only if $w_i^* = \omega_i$ for all $i$
    \item \textbf{Price dispersion}: $\KL{w^*}{\omega} > 0$ whenever prices are heterogeneous
\end{enumerate}
\end{proposition}

\paragraph{Economic interpretation.}
The KL divergence measures the ``informational distance'' between what the consumer \emph{would} choose (weights $\omega_i$) and what they \emph{actually} choose (shares $w_i^*$) given prices. It quantifies the behavioral distortion induced by price heterogeneity.

% --------------------------------------------------------------
\subsection{Welfare Sensitivity to Substitutability}
\label{subsec:welfare-sensitivity}
% --------------------------------------------------------------

The central result connects the KL divergence to how welfare responds to changes in substitutability. Recall from the generalized means chapter that for a classical mean $\Mbar_{\tilde{\rho}}$:
\begin{equation}
\frac{d \ln \Mbar_{\tilde{\rho}}}{d\tilde{\rho}} = \frac{1}{\tilde{\rho}^2} \KL{\tilde{w}}{\omega},
\end{equation}
where $\tilde{w}$ are the induced weights.

Applying this to the price index $\Pstar = \Mbar_{1-\varepsilon}(\mathbf{p}; \bom)$:

\begin{proposition}[Price Index Sensitivity]
\label{prop:price-sensitivity}
The sensitivity of the log price index to the power parameter is:
\begin{equation}
\frac{d \ln \Pstar}{d(1-\varepsilon)} = \frac{\KL{w^*}{\omega}}{(1-\varepsilon)^2}.
\end{equation}
\end{proposition}

Since indirect utility is $V = W/\Pstar$, we have $\ln V = \ln W - \ln \Pstar$. Differentiating:

\begin{proposition}[Welfare Sensitivity]
\label{prop:welfare-sensitivity}
The sensitivity of log welfare to the substitution parameter satisfies:
\begin{equation}
\frac{d \ln V}{d\varepsilon} = \frac{\KL{w^*}{\omega}}{(1-\varepsilon)^2} > 0 \quad \text{when } \KL{w^*}{\omega} > 0.
\end{equation}
Higher substitutability (higher $\varepsilon$) increases welfare whenever prices are heterogeneous.
\end{proposition}

\begin{proof}
We have $\ln V = \ln W - \ln \Pstar$. Since $W$ is independent of $\varepsilon$:
\[
\frac{d \ln V}{d\varepsilon} = -\frac{d \ln \Pstar}{d\varepsilon} = -\frac{d \ln \Pstar}{d(1-\varepsilon)} \cdot \frac{d(1-\varepsilon)}{d\varepsilon} = -\frac{\KL{w^*}{\omega}}{(1-\varepsilon)^2} \cdot (-1) = \frac{\KL{w^*}{\omega}}{(1-\varepsilon)^2}.
\]
\end{proof}

\begin{calloutimportant}[The Value of Flexibility]
This result has a striking interpretation. The \textbf{marginal value of substitutability} is proportional to the KL divergence:
\begin{itemize}[nosep, leftmargin=1.5em]
    \item When prices are equal, $\KL{w^*}{\omega} = 0$, and substitutability has no value---there is nothing to substitute toward.
    \item When prices differ, $\KL{w^*}{\omega} > 0$, and higher $\varepsilon$ allows the consumer to shift toward cheaper goods, raising welfare.
    \item The \emph{larger} the price dispersion (measured by KL divergence), the \emph{more valuable} is the ability to substitute.
\end{itemize}
\end{calloutimportant}

% --------------------------------------------------------------
\subsection{The Chain of Implications}
\label{subsec:chain}
% --------------------------------------------------------------

We can now trace the complete logic from price heterogeneity to welfare:

\begin{enumerate}[leftmargin=2em]
    \item \textbf{Price heterogeneity} ($p_i \neq p_j$) induces budget share distortion ($w_i^* \neq \omega_i$).
    
    \item \textbf{Budget share distortion} is measured by $\KL{w^*}{\omega} > 0$.
    
    \item \textbf{KL divergence} governs the sensitivity of $\ln \Pstar$ to the power parameter.
    
    \item \textbf{Price index sensitivity} translates to welfare sensitivity via $V = W/\Pstar$.
    
    \item \textbf{Higher substitutability} (higher $\varepsilon$) lowers $\Pstar$ and raises welfare.
\end{enumerate}

\begin{table}[H]
\centering
\caption{Price Dispersion and Welfare}
\label{tab:kl-welfare}
\renewcommand{\arraystretch}{1.4}
\begin{tabular}{@{}lcc@{}}
\toprule
\textbf{Condition} & \textbf{KL Divergence} & \textbf{Welfare Implication} \\
\midrule
Equal prices & $\KL{w^*}{\omega} = 0$ & $\Pstar$ invariant to $\varepsilon$; no gain from substitution \\[2pt]
Unequal prices & $\KL{w^*}{\omega} > 0$ & Higher $\varepsilon$ lowers $\Pstar$, raises $V$ \\[2pt]
Large dispersion & Large $\KL{w^*}{\omega}$ & Large welfare gains from substitutability \\
\bottomrule
\end{tabular}
\end{table}

% --------------------------------------------------------------
\subsection{Why KL Divergence?}
\label{subsec:why-kl}
% --------------------------------------------------------------

One might ask why the KL divergence---rather than variance or some other measure---is the ``right'' metric for price dispersion.

\paragraph{Structural emergence.}
The KL divergence is not an ad hoc choice. It emerges naturally from the mathematics of generalized means. The sensitivity formula $d \ln \Mbar / d\tilde{\rho} = \KL{\tilde{w}}{\omega}/\tilde{\rho}^2$ is a theorem, not an assumption.

\paragraph{Economic meaning.}
The induced weights $w_i^*$ have direct economic interpretation as budget shares. The KL divergence $\KL{w^*}{\omega}$ measures the distance between \emph{behavior} (what the consumer does) and \emph{preferences} (what the consumer would do absent price distortion).

\paragraph{Correct limiting behavior.}
The KL divergence equals zero precisely when there is no distortion (equal prices), and is strictly positive otherwise. This matches economic intuition about when substitutability matters.

\paragraph{Contrast with variance.}
Price variance, by contrast:
\begin{itemize}[nosep, leftmargin=1.5em]
    \item Does not account for preference weights $\omega_i$
    \item Depends on arbitrary price scaling
    \item Does not emerge from the optimization structure
    \item Does not directly govern welfare sensitivity
\end{itemize}

% --------------------------------------------------------------
\subsection{Application: Gains from Price Convergence}
\label{subsec:price-convergence}
% --------------------------------------------------------------

Consider an economy where prices converge toward equality---for example, through trade liberalization or technological change that reduces price dispersion.

\paragraph{Initial state.}
With heterogeneous prices $\mathbf{p}^0$, the consumer faces KL divergence $\KL{w^0}{\omega} > 0$ and achieves welfare $V^0 = W/\Pstar(\mathbf{p}^0)$.

\paragraph{Final state.}
With equal prices $p^1 = \bar{p}$ (say, the geometric mean of initial prices), the consumer achieves $\KL{w^1}{\omega} = 0$ and welfare $V^1 = W/\bar{p}$.

\paragraph{Welfare gain.}
The welfare gain from price convergence is:
\begin{equation}
\ln V^1 - \ln V^0 = \ln \Pstar(\mathbf{p}^0) - \ln \bar{p}.
\end{equation}

For small price dispersion, a Taylor expansion gives:
\begin{equation}
\ln V^1 - \ln V^0 \approx \frac{\varepsilon - 1}{2} \cdot \text{Var}_\omega(\ln p),
\end{equation}
where $\text{Var}_\omega(\ln p) = \sum_i \omega_i (\ln p_i - \overline{\ln p})^2$ is the variance of log prices weighted by $\omega_i$.

\begin{calloutnote}[Connection to Trade Theory]
This approximation underlies calculations of ``gains from trade'' in international economics. Price convergence across countries---driven by trade---generates welfare gains proportional to the variance of log prices, with the elasticity of substitution $\varepsilon$ governing the magnitude. High substitutability amplifies the gains.
\end{calloutnote}

% --------------------------------------------------------------
\subsection{Summary}
\label{subsec:kl-summary}
% --------------------------------------------------------------

The KL divergence provides a unified measure of price distortion and its welfare consequences:

\begin{enumerate}[leftmargin=2em]
    \item Budget shares $w_i^*$ are the induced weights in the price index $\Pstar$.
    
    \item The KL divergence $\KL{w^*}{\omega}$ measures how far behavior deviates from undistorted preferences.
    
    \item Zero divergence characterizes the equal-price benchmark where substitutability is irrelevant.
    
    \item Positive divergence implies that higher substitutability raises welfare: $dV/d\varepsilon > 0$.
    
    \item The magnitude of welfare sensitivity is proportional to the KL divergence itself.
\end{enumerate}

\begin{callouttip}[The Deeper Connection]
The information-theoretic interpretation runs deeper than this chapter develops. In the chapter on generalized means, we showed that the CES aggregator can be derived from entropy-regularized optimization---the Gibbs variational principle. The parameter $\rho$ governs the cost of deviating from default weights. The KL divergence that appears here is the same object that appears there, reflecting a fundamental unity between economic optimization and information theory.
\end{callouttip}
